\chapter{系统实现与关键模块}

\section{关键实现文件与模块关系}
代码结构按照“前端交互、后端服务、算法脚本分离”的思路组织。后端入口文件为 \texttt{backend/app/main.py}，负责创建 FastAPI 应用、配置跨域访问、挂载 API 路由和静态媒体路由；主要业务接口位于 \texttt{backend/app/api/routes.py}。前端入口文件为 \texttt{frontend/src/main.tsx}，核心页面由 \texttt{frontend/src/app/App.tsx} 组织。MRI 数据预处理由 \texttt{src/preprocess\_mri.py} 实现，二维生成训练与推理复用 StyleGAN2-ADA 工程，点云初始化由独立脚本完成。这样的划分让浏览器界面、后端接口、算法计算和实验产物保持相对独立。

\begin{table}[H]
\centering
\caption{系统关键实现文件及功能说明}
\begin{tabular}{>{\raggedright\arraybackslash}p{4cm}p{8.5cm}}
\toprule
文件或目录 & 主要作用 \\
\midrule
\thesispath{backend/app/main.py} & 创建 FastAPI 应用，配置 CORS，挂载 API、WebSocket 和静态媒体路由。 \\
\thesispath{backend/app/api/routes.py} & 定义系统概览、任务、生成、资产、影像分析、点云统计、报告和 AI 助手接口。 \\
\thesispath{frontend/src/app/App.tsx} & 组织 React 页面，包括仪表盘、生成工作台、影像查看器、影像分析报告、可信性评估、任务中心、三维展示和教学中心。 \\
\thesispath{frontend/src/api/client.ts} & 封装前端对 FastAPI 接口的访问，并定义图像、任务、报告和体数据等响应类型。 \\
\thesispath{src/preprocess_mri.py} & 读取 NIfTI 格式 MRI 体数据，提取轴向切片，完成归一化、缩放、低质量切片过滤和 PNG 保存。 \\
\thesispath{src/stylegan2} & 提供 StyleGAN2-ADA 的数据打包、训练、生成器加载和图像推理能力。 \\
\thesispath{src/init_pcd.py} & 将连续二维切片序列转换为 ASCII PLY 点云，作为 3DGS 三维表达的输入准备。 \\
\thesispath{src/gaussian_splatting} & 保存 3D Gaussian Splatting 框架代码，支撑 MRI 切片序列的三维高斯训练、渲染和点云导出。 \\
\thesispath{results/training_runs} & 保存训练配置、训练日志、真实样本网格和模型权重快照。 \\
\thesispath{results/generated_slices} & 保存平台根据训练权重生成的 MRI 切片结果。 \\
\thesispath{results/3dgs_runs} & 保存 3DGS 训练配置、相机参数、checkpoint、渲染图像和训练后点云结果。 \\
\bottomrule
\end{tabular}
\end{table}

\section{MRI 数据预处理模块实现}
MRI 数据预处理模块由 \texttt{src/preprocess\_mri.py} 实现。脚本接收输入目录、输出目录和目标尺寸等参数，并使用 \texttt{nibabel} 读取 \texttt{.nii} 或 \texttt{.nii.gz} 格式的 NIfTI 文件。脑部 MRI 体数据两端常包含较大背景或低信息区域，因此系统默认沿轴向提取中间约 15\% 到 85\% 的切片，以减少头顶部、颈部和空白切片对训练集的影响。

灰度处理阶段采用 2\% 到 98\% 百分位截断，而不是直接使用全局最小值和最大值进行拉伸。这种处理可以降低异常亮点、噪声和局部极值的影响。归一化后的切片被转换为 \texttt{uint8} 类型，再通过双三次插值缩放到目标分辨率。为了控制训练集质量，脚本还会统计切片平均灰度、标准差和非零像素比例，只有满足信息量要求的切片才会保存。

预处理完成后，平台可以调用 StyleGAN2 的 \texttt{dataset\_tool.py}，将 PNG 切片目录打包为训练数据。本次训练使用的数据集文件为 \texttt{mri128.zip}，训练集最大规模为 60815 张，图像分辨率为 $128\times128$。至此，医学体数据被转换为二维生成模型可以使用的训练样本。

\section{StyleGAN2 训练与切片生成实现}
StyleGAN2 训练由 FastAPI 后端服务封装。平台会在训练数据目录中查找 zip 数据包，并优先使用本项目配置的 \texttt{mri128.zip}。训练任务启动后，后端不是在接口线程中直接等待训练结束，而是调用 StyleGAN2 训练脚本启动后台子进程。子进程会接收输出目录、数据路径、GPU 数量、训练 kimg、batch size、学习率和快照保存频率等参数。这样做可以避免 React 页面长时间无响应。

训练任务启动后，系统会记录任务状态、日志路径和输出目录。前端通过任务接口或日志接口查看进度，用户不需要打开训练目录手动查找过程信息。当前项目主要使用单 GPU 训练配置，未实现多 GPU 调度。对本课题来说，单 GPU 已经能够完成 MRI 切片生成模型训练，也更便于在毕业设计环境中复现。

切片生成模块使用训练权重中的 \texttt{G\_ema} 生成器进行推理。平台可以自动查找最新网络快照，用户也可以手动指定权重路径。推理阶段，系统根据随机种子生成潜变量，并结合截断参数 \texttt{truncation\_psi} 调节输出的稳定性和多样性。生成器输出张量从 $[-1,1]$ 转换到 $[0,255]$ 后保存为 PNG 图像。目前项目已经形成编号为 010000 的模型快照，并生成了示例切片。

该模块的重点是把训练产物真正接入平台。用户通过界面设置生成数量、随机种子和截断参数后，前端向后端生成接口发送请求。后端加载权重、执行推理、保存图像，并把输出路径和图像资产信息返回给前端。这样设计后，论文插图整理更方便，答辩时也能清楚说明模型输入、参数和输出之间的关系。

\section{连续切片序列与点云初始化实现}
连续切片序列由后端生成服务中的 \texttt{generate\_sequence} 函数生成。系统分别根据两个随机种子构造潜变量 $z_a$ 和 $z_b$，再对二者进行线性插值，得到一组连续变化的潜变量。每个插值潜变量输入 StyleGAN2 生成器后输出一张 MRI 切片，并按 \texttt{slice\_0000.png}、\texttt{slice\_0001.png} 等规则保存，为 2D/3D 对比展示提供连续输入。

点云初始化由 \texttt{init\_pcd.py} 完成。脚本读取连续切片目录中的 PNG 图像，并把灰度值大于阈值 10 的像素转换为空间点。对于第 $z$ 张切片中的像素坐标 $(x,y)$，系统将其映射为三维点 $(x,y,z\cdot spacing)$，其中 \texttt{spacing} 表示切片间距。像素灰度值归一化后写入 RGB 颜色字段，所有点最终以 ASCII PLY 格式保存。本项目已生成 \texttt{results/3d\_model\_test/init\_point\_cloud.ply}，该文件包含 477426 个顶点，大小约 13 MB。

通过该模块，二维切片被整理为三维点云。每个点都包含空间坐标和灰度颜色。点云既可以作为 3DGS 训练输入，也可以展示二维生成结果如何过渡到三维表达。

\section{3DGS 三维重建模块实现}
3DGS 模块围绕 MRI 切片序列、初始点云、相机参数、训练日志和输出点云设计统一接口。该实现路线对应本文的工程目标：先把二维生成结果组织为三维输入，再通过 3DGS 优化形成高斯表达。这里的关键不是重新提出 3DGS 算法，而是把 MRI 切片整理成 3DGS 能够读取和训练的工程输入。

3DGS 模块输入包括三类数据。连续 MRI 切片序列提供渲染监督，也就是训练时需要拟合的目标图像；根据切片灰度和空间位置初始化得到的 PLY 点云提供空间起点；相机参数文件用于说明每张切片对应的投影关系。普通 3DGS 场景通常借助 COLMAP 从多视角照片中估计相机位姿，但 MRI 切片并没有真实相机拍摄过程。基于这一差异，系统使用切片索引、像素分辨率和层间间隔构建规则相机参数。当前训练目录中的 \texttt{cameras.json} 包含 180 个相机条目，图像分辨率为 $256\times256$，焦距参数 $fx$ 与 $fy$ 均为 426.67。

本文没有直接使用 NIfTI affine 矩阵生成真实物理空间坐标，而是采用规则化切片坐标完成工程验证。这样处理的优点是实现简单、输入稳定，也便于与生成切片序列衔接；不足是不能完整保留真实扫描中的空间方向和体素间距。后续如果要提升医学空间一致性，需要进一步把 affine 矩阵、真实层厚和体素间距纳入相机参数构建过程。

训练实现上，系统使用 180 张连续 MRI 切片和 300000 个初始化点构建 3DGS 输入数据，并把结果保存到 3DGS 主训练目录。训练过程记录相机参数、配置参数、训练日志、checkpoint、渲染图像、参考切片和迭代点云。第 30000 轮输出的完整 PLY 文件头部显示包含 556424 个高斯点，同时保存的 \texttt{chkpnt30000.pth} 用于记录训练后的模型状态。为了兼顾浏览器加载速度和交互流畅性，系统还从完整点云中导出 60000 点预览 PLY，用于前端点云展示。

从平台使用角度看，三维重建模块接收切片目录和输出目录等核心参数，并返回训练日志、点云路径和渲染结果路径。系统先完成点云初始化，再组织训练迭代次数、输出目录、相机参数和渲染结果目录等信息。用户可以在同一界面中查看、下载和归档三维结果，不需要在多个脚本和文件夹之间反复切换。

这一部分被归入 3DGS 三维重建模块，是因为项目已经形成真实切片序列、PLY 点云文件、训练日志、checkpoint 和渲染截图。模块完成了从二维切片组织到三维高斯表达结果的工程闭环。第 6 章将结合训练指标和渲染效果继续分析。

\section{FastAPI 后端与 React 前端实现}
后端基于 FastAPI 实现。入口文件 \texttt{backend/app/main.py} 创建应用实例，并挂载 \texttt{/api} 前缀下的业务路由。系统使用 CORS 中间件允许前端访问后端接口。API 路由文件 \texttt{backend/app/api/routes.py} 定义系统概览、任务查询、预处理任务、训练任务、切片生成、连续序列生成、图像资产、模型资产、点云资产、影像分析、影像分析报告、点云统计和 AI 助手等接口。对于需要持续更新的任务状态，后端还提供 \texttt{/ws/tasks} 和 \texttt{/ws/logs/\{job\_id\}} 等 WebSocket 接口。

前后端通信采用比较直接的请求流程。用户在 React 页面中填写参数并点击按钮后，前端通过 \texttt{client.ts} 调用 FastAPI 接口。后端校验参数、调用对应脚本或读取结果文件，再把 JSON 数据、图像路径、点云路径或任务状态返回给前端。对于训练这类耗时任务，接口通常先返回任务编号，后续状态由任务查询接口和 WebSocket 日志接口更新。这个流程把“提交任务”和“等待结果”分开，能够减少页面阻塞。

前端基于 React 和 Vite 实现。项目入口为 \texttt{frontend/src/main.tsx}，核心应用组件为 \texttt{frontend/src/app/App.tsx}。前端页面包括仪表盘、生成工作台、影像查看器、影像分析报告、可信性评估、任务中心、三维展示和教学中心等模块。前端通过 \texttt{frontend/src/api/client.ts} 统一访问后端接口。该客户端封装概览、切片生成、图像资产、影像查看、图像分析和 AI 助手等请求，并将后端返回结果转换为页面可以直接使用的数据结构。

影像展示部分由多个前端查看器组件完成。医学切片浏览、点云查看和体渲染展示分别对应切片查看器、点云查看器和体数据查看器。平台还使用图表、质量与风险提示面板和可信性证据面板展示生成结果的质量、结构和风险提示。启动脚本 \texttt{run\_fastapi\_react.bat} 会分别启动后端 Uvicorn 服务和前端 Vite 服务，本地后端默认运行在 \texttt{127.0.0.1:8000}，前端默认运行在 \texttt{127.0.0.1:5180}。

\section{AI 助手功能实现}
AI 助手以悬浮按钮和聊天面板形式集成在 React 前端中。系统设置了面向 MRI 医学图像拟生成平台的提示词。助手回答时主要围绕 StyleGAN2、3DGS、MRI 图像处理和平台使用流程展开。前端把用户问题、当前图像、分析结果和近期任务状态传给后端，后端通过 \texttt{/api/assistant/chat} 返回回复。当前实现包含本地规则解释逻辑，并保留 Bocha 搜索接口，用于辅助教学说明。

AI 助手的定位是教学辅助和参数解释，不是核心算法依赖。用户可以询问截断参数、随机种子、点云初始化、3DGS 三维表达和 MRI 切片生成等问题。AI 助手依赖外部接口，可用性会受到网络和 API 服务影响。为保证平台稳定运行，数据处理、模型训练、图像生成和点云初始化功能都保持独立。即使 AI 助手不可用，平台的主要毕业设计功能仍能正常完成。
